\documentclass[10pt,twocolumn]{article}
\usepackage[T1]{fontenc}
\usepackage[utf8]{inputenc}
\usepackage{lmodern}
\usepackage[margin=1.8cm,top=2.4cm,bottom=2cm,columnsep=0.6cm]{geometry}
\usepackage{fancyhdr}
\usepackage{titlesec}
\usepackage{booktabs}
\usepackage{amsmath,amssymb,amsfonts}
\usepackage{microtype}
\usepackage{xcolor}
\usepackage{mdframed}
\usepackage{parskip}
\usepackage{enumitem}
\usepackage{hyperref}

\definecolor{rulered}{RGB}{180,30,30}
\definecolor{boxgray}{RGB}{245,245,245}
\definecolor{keywordgray}{RGB}{100,100,100}

\pagestyle{fancy}
\fancyhf{}
\fancyhead[L]{\textsc{\small Claude Mag}}
\fancyhead[C]{\small\textit{Machine Learning Research Digest}}
\fancyhead[R]{\small 2026-11-04}
\fancyfoot[C]{\small\thepage}
\renewcommand{\headrulewidth}{0.8pt}

\titleformat{\section}{\normalsize\bfseries\color{rulered}}{}{0pt}{}%
  [\vspace{-4pt}\color{rulered}\rule{\columnwidth}{0.4pt}]
\titlespacing{\section}{0pt}{8pt}{4pt}

\hypersetup{colorlinks=true,urlcolor=rulered,linkcolor=black}

\begin{document}
\setlength{\parindent}{0pt}
\setlength{\parskip}{4pt}

{\small\color{keywordgray}\textbf{Keywords:}
  multi-agent reasoning \textbullet{}
  strategy optimization \textbullet{}
  meta-learning \textbullet{}
  agentic AI}\par\vspace{4pt}

{\LARGE\bfseries\raggedright
  Strategist: Designing Agentic Reasoning at Scale}\par\vspace{4pt}

{\large\itshape\raggedright
  A Meta-Agent Builds a Growing Library of Executable Strategy
  Blueprints, Winning 24 of 30 Benchmarks at 48\% Lower Cost}\par\vspace{6pt}

{\small\textbf{By} Taiming Lu, Qi Xu, Zhigang Hua, Zhuang Liu}\par\vspace{2pt}

{\small\color{keywordgray}arXiv:2611.03618 \textbullet{}
  cs.AI / cs.CL \textbullet{} NeurIPS 2026 Poster}
\par\vspace{2pt}\noindent{\color{rulered}\rule{\columnwidth}{0.6pt}}\vspace{6pt}

Strategist is a meta-agent that automatically designs inference strategies
for new tasks by representing each strategy as executable Python code,
building a reusable library of blueprints that grows across tasks.
Evaluated on 30 reasoning benchmarks, Strategist wins 24 and achieves
8.5 absolute percentage points higher average accuracy than the strongest
fixed-strategy baseline at 48\% lower inference cost. Strategies transfer
across model backends the meta-agent never optimized on.

\begin{mdframed}[backgroundcolor=boxgray,linecolor=rulered,linewidth=1pt,
    innerleftmargin=6pt,innerrightmargin=6pt,
    innertopmargin=6pt,innerbottommargin=6pt]
\textbf{\small AT A GLANCE}\par\vspace{4pt}
\begin{description}[leftmargin=0pt,labelsep=4pt,itemsep=2pt,parsep=0pt,
    font=\small\bfseries]
  \item[Problem:] Fixed pipeline templates (chain-of-thought, tree-of-thought,
    debate) do not generalize across domains and require manual design.
  \item[Why it matters:] The topology of multi-agent reasoning is a critical
    variable for performance; automating its design could yield consistent
    improvements across tasks without human effort.
  \item[Method:] Strategies are Python programs (blueprints); a meta-agent
    retrieves related blueprints, proposes candidates, evaluates them on a
    dev set, revises them, and stores all valid candidates in a library.
  \item[Result:] 8.5pp accuracy improvement and 48\% cost reduction over
    the best fixed strategy, on 24 of 30 benchmarks.
  \item[Evidence:] 30 benchmarks spanning math, coding, science, and
    commonsense; transfer experiments to Gemini-2.5 backend.
\end{description}
\end{mdframed}

\section{The Big Picture}

Multi-agent systems for reasoning typically rely on templates: chain-of-thought
prompts the model to reason step-by-step, tree-of-thought expands multiple
reasoning branches, and debate pits agents against one another. These
structures are designed by researchers for specific task types and do not
automatically adapt to new domains. Meanwhile, automated prompt engineering
optimizes over discrete tokens but leaves the structural topology of the
reasoning process fixed.

The key insight of Strategist is that the topology itself---how agents
are orchestrated, how outputs are routed and verified---is a programmable
object that can be optimized. Representing strategies as Python programs
makes them composable, executable, and amenable to code-level search by
a language model.

\section{Why Existing Approaches Fall Short}

\textbf{Fixed templates} (chain-of-thought, majority voting) are designed
for average-case performance across broad task categories. They cannot
capture the structural insights that are most effective for a specific
task type, such as using a symbolic calculator for arithmetic sub-steps
or running parallel hypotheses for spatial reasoning.

\textbf{Prompt optimization} methods (AutoPrompt, DSPy, APE) treat the
strategy as a fixed function and optimize over its textual inputs, but
do not change the number of agents, the verification steps, or the
routing logic. They are limited to the structural assumptions baked into
the initial template.

\textbf{Human-designed multi-agent systems} require expert knowledge and
significant engineering time, and the resulting designs are rarely
reusable across domains.

\section{Core Mechanism: Strategy as Code}

Each strategy in Strategist is a Python program, or \emph{blueprint},
that takes a task description and input as arguments and returns an
answer by orchestrating one or more LLM calls. A blueprint may:
\begin{itemize}[itemsep=2pt,parsep=0pt]
  \item Decompose the problem into typed sub-problems and dispatch each
    to a specialized agent.
  \item Route outputs conditionally based on model confidence or
    consistency checks.
  \item Call helper functions from other library blueprints as
    sub-routines.
  \item Aggregate results via voting, scoring, or symbolic verification.
\end{itemize}

This representation makes strategies fully executable and inspectable,
enables compositional reuse, and gives the meta-agent a structured
search space that is richer than prompt tokens alone.

\section{The Design Loop}

Given a new task $T$ and a small development set $\mathcal{D}$:

\begin{enumerate}[itemsep=2pt,parsep=0pt]
  \item \textbf{Retrieve:} Encode $T$ with a frozen text encoder and
    retrieve top-$k$ blueprints from the library by cosine similarity.
  \item \textbf{Propose:} The meta-agent generates 3--5 candidate
    blueprints as Python code, guided by retrieved examples and the
    target task description.
  \item \textbf{Evaluate:} Execute each candidate on $\mathcal{D}$;
    record accuracy $a_i$ and total token cost $c_i$.
  \item \textbf{Revise:} The meta-agent reviews evaluation results and
    produces a final revised blueprint for $T$.
  \item \textbf{Store:} All candidates and the final blueprint are added
    to the library with metadata $(a_i, c_i, \text{task-type tags})$.
\end{enumerate}

The loop runs for 3 rounds per task, with early stopping if the best
candidate exceeds a target accuracy of $\alpha^* = 0.9$.

\section{The Strategy Library}

The library is initialized with 20 hand-designed baselines. After
processing all 30 benchmarks in order, the library contains
$\sim$150 blueprints. The compositional reuse mechanism means that
blueprints are not monolithic: a later task (e.g., GPQA) can invoke
the ``verify-arithmetic'' sub-routine discovered for GSM8K and the
``multi-perspective synthesis'' routine from debate-style blueprints.

\section{What the Guarantee Says}

Strategist provides no PAC-learning or regret guarantee. The empirical
claim is competitive: on the 30-benchmark suite, the designed strategies
outperform all fixed baselines (including the best oracle fixed strategy
per task). No theoretical bound on the number of meta-agent iterations
needed to find a good strategy is given. The transfer result
(Section~7) is empirical and suggests cross-model generality, but
the mechanism is not formally analyzed.

\section{Experimental Findings}

\begin{itemize}[itemsep=2pt,parsep=0pt]
  \item \textbf{Accuracy (30 benchmarks):} Strategist wins 24 of 30;
    average improvement is \textbf{8.5pp} over the strongest fixed
    baseline (which is the oracle best-fixed-strategy per task).
  \item \textbf{Cost:} Despite higher quality, Strategist uses
    \textbf{48\% fewer tokens} on average, because it avoids
    unnecessary LLM calls and uses cheaper model tiers for sub-tasks.
  \item \textbf{Library growth:} Average accuracy on a held-out set of
    5 benchmarks improves by 3.1pp after processing the first 25 tasks,
    confirming compositional library benefit.
  \item \textbf{Transfer to Gemini-2.5:} Strategies discovered with
    GPT-6 improve Gemini-2.5 chain-of-thought accuracy by \textbf{6.1pp}
    on average, confirming cross-model transferability.
\end{itemize}

\section{Ablations and Interpretation}

\textbf{Without retrieval:} Removing library retrieval (proposing
strategies from scratch) reduces accuracy by 4.4pp, confirming that
the library enables meaningful compositional reuse.

\textbf{Strategies as text vs.\ code:} Replacing Python blueprints with
natural-language strategy descriptions reduces accuracy by 2.9pp, as
natural-language strategies are less precise and cannot directly invoke
verified sub-routines.

\textbf{Without iterative revision:} Single-round (no revision step)
performance is 3.7pp below the full 3-round loop, with diminishing
returns in rounds 2 and 3.

\textbf{Meta-agent backbone:} Replacing GPT-6 as the meta-agent with
GPT-5 reduces strategy quality by 2.1pp on average, but the benefit of
strategy design over no design is maintained across all backbones tested.

\section*{Reference}

Taiming Lu, Qi Xu, Zhigang Hua, Zhuang Liu.
\textbf{Strategist: Designing Agentic Reasoning at Scale.}
arXiv:2611.03618. NeurIPS 2026.
\url{https://arxiv.org/abs/2611.03618}

\end{document}
